\section{Orden verbal con modales}
\textbf{Regla base.} El modal ocupa la segunda posición; el infinitivo principal va al final.

\textbf{Patrones}
\begin{itemize}[leftmargin=*]
  \item Declarativa: [Tiempo/Lugar] + [Modal (V2)] + [Sujeto si no abre] + \ldots{} + [Infinitivo].
  \item Sí/No: [Modal] + [Sujeto] + \ldots{} + [Infinitivo]? \ \textit{Ej.: Kun jij helpen?}
  \item Negación: \textit{niet} se coloca antes del infinitivo final.\ \textit{Ej.: Wij willen niet vertrekken}.
\end{itemize}

\textbf{Ejemplos}
\begin{itemize}[leftmargin=*]
  \item \textit{Morgen kan ik niet komen}. (Mañana no puedo venir.)
  \item \textit{In Antwerpen wil zij wonen}. (Ella quiere vivir en Amberes.)
  \item \textit{Moeten jullie om acht uur beginnen?}
\end{itemize}

\textbf{Mini práctica}
\begin{itemize}[leftmargin=*]
  \item Reordena con tiempo inicial: \textit{ik / kunnen / vandaag / niet / fietsen}.
  \item Forma una pregunta de sí/no con \textit{willen}.
\end{itemize}
